La enseñanza de la programación en los cursos introductorios, comúnmente denominados CS1 (\textit{Computer Science 1}), representa uno de los desafíos pedagógicos más críticos en la educación superior contemporánea \cite{alshaikh2020socratic}. Históricamente, los programas académicos de informática experimentan altas tasas de abandono en etapas iniciales, con porcentajes de deserción que frecuentemente oscilan entre el 30\% y el 40\%. Esta problemática no responde únicamente a una falta de interés, sino a la complejidad de asimilar simultáneamente tanto las reglas gramaticales del código como la lógica necesaria para resolver problemas.

Para un estudiante que se enfrenta por primera vez a la programación, el entorno de desarrollo puede resultar abrumador. La combinación de conceptos abstractos y la rigidez de las herramientas genera una frustración profunda. Al no contar todavía con un modelo mental sólido, los alumnos tienen grandes dificultades para entender por qué su código no funciona, sintiéndose a menudo incapaces de avanzar de forma autónoma. Idealmente, un profesor debería estar disponible para ofrecer pistas en el momento justo, pero en las facultades actuales, con clases masificadas, es físicamente imposible atender a cada estudiante de forma individualizada. Esta limitación es la que hace necesario buscar sistemas de tutoría automatizados que ofrezcan ese apoyo constante en tiempo real.